\subsection*{1\quad 日内新信息怎样改变购电决策}

问题二只能在每天零点制定一次购电计划，白天即使发现光伏偏离预期，也不能修改普通订单。问题三允许在6:00、12:00和18:00收到新光伏预报后，重新安排尚未执行的部分。因此，本问的关键不只是把预报做得更准，还要判断修改订单能否抵消调整费，并让电池在新旧计划之间平稳衔接。

我们把一天划分为144个10分钟区间，记日期为 $d$，日内区间为 $k=1,\ldots,144$，每段长度为 $\Delta t=1/6$ 小时。第1段对应0:00--0:10，第144段对应23:50--24:00。零点形成当天原计划，6:00、12:00和18:00分别在第37、73和109段开始前更新；已经过去的区间保持不变，只调整未来订单。附件2中的实际负载和实际光伏只在相应时段结束后成为历史信息，附件3才是各发布时间能够看到的未来光伏预报。

新预报并不意味着每次改单都划算。如果预报只发生小幅变化，节省的购电量可能还不够支付调整费；如果电池已经存有足够电量，也可能无需改约。为选出与这种信息结构相符的方案，我们设置四个候选：直接使用官方预报；用历史误差融合官方预报；用提升树修正官方预报；在融合预报上再加入近期误差准备。四个候选都接入相同的日内滚动调度，因而能够把预报方法的差别落实到最终电费上。

官方预报给出了当天最直接的信息，但它按小时发布且仍有误差；近期同刻历史能补充稳定规律，却不能替代当天预报；提升树可以表达非线性，但较难直接说明每次修正来自哪里。另一方面，少估缺口会触发五倍紧急电费，所以模型还必须专门处理低估风险。基于这些特点，我们预先选择“历史融合＋动态准备＋滚动规划”作为主方案，再用全年费用检验这一判断。最终模型命名为：

\textbf{历史融合预报--动态安全准备--日内滚动规划模型。}

\subsection*{2\quad 用历史记录修正光伏预报}

负载和光伏的形成规律不同，因此仍然分开估计。负载沿用问题二的同星期历史参考：以过去四个相同星期的同一时刻均值作为基础预测。光伏则先读取附件3在当前时点发布的预报，再用此前已经兑现的实际光伏判断这份预报应当相信多少。

附件3按整点给出预测，而购电计划按10分钟填写。对同一发布时间的相邻两个整点预测采用线性插值，将其转换为10分钟功率。该处理保留了整点值和相邻时刻的变化方向，只是把小时信息平滑分配到10分钟尺度，并不把插值结果当作新获得的真实观测。

记当前发布时间为 $r$，附件3给出的光伏预测为 $F_{r,t}$，最近七个可用同刻实际值形成的历史参考为 $H_{r,t}$。将未来24小时按提前1--6、7--12、13--18和19--24小时分组，在每一组中用过去28天、且在发布时间 $r$ 前已经能够核验的样本确定融合权重：
\begin{equation}
\widehat V_{r,t}=w_{r,g}F_{r,t}+(1-w_{r,g})H_{r,t},
\qquad 0\le w_{r,g}\le1,
\label{eq:q3-fusion}
\end{equation}
\begin{equation}
w_{r,g}=\operatorname{clip}_{[0,1]}
\frac{\sum_{j\in\mathcal H_{r,g}}(F_j-H_j)(V_j-H_j)}
{\sum_{j\in\mathcal H_{r,g}}(F_j-H_j)^2}.
\label{eq:q3-weight}
\end{equation}
其中，$\mathcal H_{r,g}$ 只包含当时已经发生的历史记录。式\eqref{eq:q3-weight}选择使历史平方误差最小的权重：官方预报过去更可靠时，$w$ 自动增大；历史同刻曲线更可靠时，$w$ 自动减小。它会随发布时间和提前时长变化，并不是人为指定一个固定比例，也不是用熵权法分配权重。

\subsection*{3\quad 根据近期净需求误差留出准备量}

只提高光伏预测精度仍不能解决风险不对称的问题。少估1 kWh缺口，运行时可能要按五倍电价紧急补购；多准备1 kWh普通电，也可能因为负载较低或电池已满而无法利用。因此，我们不直接把订单整体放大，而是从近期净需求误差中估计每个时段应保留多少安全余量。

对过去已经兑现的记录，定义10分钟净需求误差为
\begin{equation}
\varepsilon_{i,k}=
\big[(L_{i,k}-V_{i,k})-(\widehat L_{i,k}-\widehat V_{i,k})\big]\Delta t.
\label{eq:q3-error}
\end{equation}
正值表示当时把缺电量看小了，负值表示计划偏保守。每次更新时，只使用此前28天的误差，并把一天分为8个3小时组。同组时段面对相近的光照和用电阶段，样本数量又足以支撑分位估计。

近期误差比很久以前的误差更能反映当前天气和用电水平。若样本 $j$ 距当前完整相隔 $a_j$ 天，赋予权重
\begin{equation}
\omega_j=\frac{\beta^{a_j}}{\sum_l\beta^{a_l}},\qquad \beta=0.7.
\label{eq:q3-decay}
\end{equation}
这表示最近一天的相对权重为1，再早一天为0.7，再早两天为0.49。其半衰期约为1.94天，既保留数周的可用记录，又让近期变化起主要作用。最后取同组误差的加权0.60分位，并把负值截为0：
\begin{equation}
m_{r,k}=\max\left\{0,Q_{0.60}^{(\omega)}(\varepsilon_{i,k})\right\}.
\label{eq:q3-margin}
\end{equation}
0.60高于中位数，是因为低估缺口的代价更高；它又没有取到过高位置，以免普通电准备过多。每次收到新预报后都重新计算并完整保留该准备量，而不是随着日内更新逐次撤掉。

\subsection*{4\quad 滚动购电与储能调度模型}

\textbf{预计供需。} 由负载预测 $\widehat L_t$、融合光伏预测 $\widehat V_t$ 和准备量 $m_t$ 得到
\begin{equation}
r_t=(\widehat L_t-\widehat V_t)\Delta t+m_t,\qquad
n_t=\max(r_t,0),\quad a_t=\max(-r_t,0).
\label{eq:q3-net}
\end{equation}
$n_t$ 表示预计仍需外网或电池补充的电量，$a_t$ 表示考虑准备量后预计能够储存的光伏富余。

\textbf{决策变量。} 令 $q_t$ 为普通购电订单，$x_t$ 和 $y_t$ 分别为外网充电量和光伏充电量，$z_t$ 为电池向负载送出的电量，$e_t$ 为规划中预留的紧急购电量，$u_t$ 为已经下单但预计无法接收的普通电量，$S_t$ 为区间结束时的内部储电量。除 $S_t$ 外，各变量均按设备外侧、每10分钟的 kWh 计量。

\textbf{供需与储能约束。} 光伏先直接满足负载，式\eqref{eq:q3-net}已经扣除了这部分直供量。其余电量满足
\begin{equation}
q_t-x_t+z_t+e_t-u_t=n_t.
\label{eq:q3-balance}
\end{equation}
普通电可以直接供负载，也可以通过 $x_t$ 给电池充电；若订单过多且电池存不下，则由 $u_t$ 记录未接收但仍需结算的部分。充、放电效率均为90\%，所以
\begin{equation}
S_t=S_{t-1}+0.9(x_t+y_t)-\frac{z_t}{0.9},qquad
1200\le S_t\le10800.
\label{eq:q3-storage}
\end{equation}
设备功率上限为5000 kW，换算到10分钟后，每段最多充入或放出 $5000\Delta t=833.33$ kWh：
\begin{equation}
\begin{aligned}
&x_t,y_t,z_t,e_t,u_t,q_t\ge0,\qquad y_t\le a_t,\\
&x_t+y_t\le5000\Delta t,\qquad z_t+e_t\le n_t.
\end{aligned}
\label{eq:q3-bounds}
\end{equation}
每天的初始库存取上一天24:00的实际剩余电量，不额外要求日首和日末相等。

\textbf{购电与调整费用。} 记零点确定的原订单为 $q_t^0$，执行前最后一次调整后的订单为 $q_t^F$。按照采用的结算口径，多个时点对同一时段的修改只看最终订单相对零点订单的净变化，不重复累计中间改单。普通结算与实际总费用分别为
\begin{equation}
C^{\rm ord}=\sum_t\left[p_tq_t^F+\frac12p_t|q_t^F-q_t^0|\right],
\label{eq:q3-ordinary}
\end{equation}
\begin{equation}
C=\sum_{d,k}\left[p_kq_{d,k}^F+\frac12p_k|q_{d,k}^F-q_{d,k}^0|
+5p_ke^{\rm actual}_{d,k}\right].
\label{eq:q3-bill}
\end{equation}
例如，原计划100 kWh、单价1元，后来减至80 kWh，则普通订单费为80元，调整费为10元，相关费用共90元；原计划的100元不再重复相加。

\textbf{目标函数。} 在每次决策时，对当前时刻至次日24:00进行规划，最小化预计普通购电、调整和紧急购电费用：
\begin{equation}
\min J=\sum_{t\in\mathcal T_r}
\left[p_tq_t+\frac12p_t\delta_t+5p_te_t\right]-\nu S_T,
\label{eq:q3-objective}
\end{equation}
\begin{equation}
\delta_t\ge q_t-q_t^0,qquad \delta_t\ge q_t^0-q_t,qquad
\nu=0.9p_{00:00}.
\label{eq:q3-absolute}
\end{equation}
$\delta_t$ 将绝对值调整费写成线性约束。末端库存项用于避免规划在边界前无节制地放空电池，$\nu$ 只代表余电未来可用价值，不作为售电收入计入式\eqref{eq:q3-bill}。

\subsection*{5\quad 模型求解}

给定某次发布时间能够获得的信息后，式\eqref{eq:q3-net}--\eqref{eq:q3-absolute}中的目标和约束都是线性的，因此采用线性规划求解。每个比赛日按以下顺序运行：
\begin{enumerate}
\item 零点读取当时已有历史、当前真实库存和当天光伏预报，将整点预报插值到10分钟，计算负载预测、融合光伏预测及误差准备量；
\item 从零点规划到次日24:00，保存当天前144段普通订单作为 $q^0$，次日部分只帮助安排跨日库存；
\item 6:00、12:00和18:00收到新预报后，以该时刻真实库存为新初值，重新计算未来预测和准备量；已经执行的订单锁定，只求解当天尚未执行的部分；
\item 执行接下来的6小时。实际光伏先供负载，剩余普通订单和光伏在功率、容量允许时充电；普通订单不足时先由电池补缺，仍不足的部分才紧急购买；
\item 24:00按最终订单相对零点订单的净变化结算，日末真实库存传给下一天。
\end{enumerate}

2月1日零点库存取1月连续运行得到的8592.69 kWh。2月1日至12月31日共334天，每天求解4次，合计1336次线性规划均得到可行最优解。最大约束残差为 $1.82\times10^{-12}$，最大对偶差为 $1.31\times10^{-10}$，说明保存的方案满足供需、库存和功率限制。

\subsection*{6\quad 从候选比较到最终策略}

前面的题意分析使我们优先选择历史融合、误差准备与滚动规划的组合，但这一预选仍需接受同条件检验。为避免只凭预测误差选模型，我们让各候选使用相同负载预测、设备参数、初始库存、实际供电规则和结算方法，只改变光伏处理和风险准备。表\ref{tab:q3-candidates}给出完整334天的实际费用。
\begin{inlinetable}
\caption{问题三候选方案的334天费用比较}\label{tab:q3-candidates}
\resizebox{\linewidth}{!}{\begin{tabular}{lrrrr}\toprule
方案 & 普通结算/万元 & 紧急购电/万元 & 总费用/万元 & 紧急购电天数\\\midrule
\tableinput{tables/candidates.tex}
\bottomrule\end{tabular}}
\end{inlinetable}

直接采用官方预报的总费用为1417.51万元。利用历史误差融合后降至1408.87万元，说明官方预报和近期同刻记录确实能够互相补充。提升树虽然可以表达更复杂的非线性关系，但用于下一阶段决策的短期误差和最终费用都没有占优。更值得注意的是，初始融合方案的普通费用较低，紧急费用却达到132.51万元，说明“平均上更准”仍可能在少数关键时段低估缺口。

在融合方案上加入近期净需求误差准备后，总费用降为1314.65万元，比初始融合方案节省94.22万元。普通结算增加17.98万元，但紧急费用减少112.20万元，减少的高价补购远大于增加的普通准备。提升树能够表达复杂关系，但本题的官方预报已经提供了较强先验，树模型的短期修正没有换来更低费用；历史融合可以清楚说明官方预报与近期记录各起什么作用，再配合准备量处理高价缺口，结构上更符合题意。由此，费用结果与建模前的判断一致，最终采用\textbf{历史融合预报--动态安全准备--日内滚动规划模型}。

\subsection*{7\quad 全年结果与指定日期分析}

最终方案在334天内的普通结算为\textbf{12943484.52元}，紧急购电费为\textbf{203050.97元}，总费用为\textbf{13146535.49元}；紧急购电量为35685.12 kWh，共有164天出现过至少一个紧急购电区间。年末内部库存为8514.41 kWh。图\ref{fig:q3-monthly}显示，11个月相对初始融合方案都实现节省，主要收益来自减少紧急购电。
\begin{inlinefigure}
\includegraphics[width=.96\linewidth]{final_monthly.pdf}
\caption{最终方案的月度费用及相对初始融合方案的节省}\label{fig:q3-monthly}
\end{inlinefigure}

月度均有节省不代表每天都更便宜。逐日比较中，250天节省、59天增加、25天基本相同；7月7日最多节省23360.46元，6月29日则增加1041.73元。模型付出的代价还包括443358.10 kWh未接收普通订单，说明安全准备能够减少高价缺口，但不能被描述为没有代价的改进。

题目指定日期的零点原计划和调整后订单见表\ref{tab:q3-original}、表\ref{tab:q3-adjusted}。正文只列六个代表时段，全天总量按完整144段求和。
\begin{inlinetable}
\caption{指定日期零点原计划的部分时段购电量（kWh）}\label{tab:q3-original}
\resizebox{\linewidth}{!}{\begin{tabular}{lrrrr}\toprule
区间 & 3月20日 & 6月21日 & 9月23日 & 12月21日\\\midrule
\tableinput{tables/original_orders.tex}
\bottomrule\end{tabular}}
\end{inlinetable}

\begin{inlinetable}
\caption{指定日期调整后订单及全天费用（电量单位：kWh）}\label{tab:q3-adjusted}
\resizebox{\linewidth}{!}{\begin{tabular}{lrrrr}\toprule
区间或项目 & 3月20日 & 6月21日 & 9月23日 & 12月21日\\\midrule
\tableinput{tables/adjusted_orders.tex}
\bottomrule\end{tabular}}
\end{inlinetable}

表\ref{tab:q3-storage}汇总各4小时内的实际充放电量。一个4小时段内同时出现充、放电，表示其中不同10分钟区间采取了不同动作，并非同一时段同时充放。四天日末库存都直接传给次日，没有人为恢复到日初水平。
\begin{inlinetable}
\caption{指定日期的实际充放电及日首尾库存（kWh）}\label{tab:q3-storage}
\resizebox{\linewidth}{!}{\begin{tabular}{llrrrr}\toprule
区间 & 项目 & 3月20日 & 6月21日 & 9月23日 & 12月21日\\\midrule
\tableinput{tables/storage.tex}
\bottomrule\end{tabular}}
\end{inlinetable}

四个指定日期都只出现少量紧急购电，具体时段见表\ref{tab:q3-emergency}。图\ref{fig:q3-days}进一步展示其中紧急量最大和最小的3月20日、9月23日。大部分普通订单后的缺口都由电池补上；3月20日20:40附近库存接近下限，仍有95.59 kWh无法由电池提供，只能紧急购入。9月23日仅在20:30出现6.38 kWh补购。可见是否需要紧急购电取决于缺口发生的具体时刻、当时库存和放电速度，不能只比较全天电量总和。
\begin{inlinetable}
\caption{指定日期的紧急购电明细}\label{tab:q3-emergency}
\begin{tabular}{llrr}\toprule
日期 & 发生区间 & 电量/kWh & 费用/元\\\midrule
\tableinput{tables/emergency.tex}
\bottomrule\end{tabular}
\end{inlinetable}

\begin{inlinefigure}
\includegraphics[width=.98\linewidth]{specified_days.pdf}
\caption{两个指定日期的供需、库存和缺口补充过程}\label{fig:q3-days}
\end{inlinefigure}

\subsection*{8\quad 日内预报带来的实际费用价值}

问题三要求评价日内预报的价值。若只把“零点全天不改订单”和“使用日内新预报”直接相减，会把实时看到实际供需、重新掌握电池库存以及允许改约的作用一并算进去。为分开这些影响，在直接采用官方预报的同一策略下设置三个逐步增加信息的对照：零点订电后全天不改；允许日内反馈和改约，但不读取后发官方预报；在此基础上再使用6:00、12:00、18:00的新预报。

如图\ref{fig:q3-value}所示，三种设置总费用依次为1614.80、1439.29和1417.51万元。开放日内反馈与改约先节省175.51万元，在此基础上使用新官方预报再节省21.78万元。前者体现“根据实际运行修正未来安排”的价值，后者才是新增官方预报在该对照顺序下带来的增量价值。两者合计197.28万元，不能全部归因于新预报。
\begin{inlinefigure}
\includegraphics[width=.96\linewidth]{information_value.pdf}
\caption{日内反馈、改约与新增官方预报的费用价值}\label{fig:q3-value}
\end{inlinefigure}

因此，对题目最后提出的“是否需要引入6:00、12:00和18:00预报”这一问题，本文的回答是：\textbf{需要引入}。在日内反馈和改单权限完全相同的条件下，使用后发预报把费用从1439.29万元降到1417.51万元，额外节省21.78万元，约为1.51\%。这说明新预报总体上有实际价值，但不表示三个发布时间同样重要，也不表示每次收到预报都必须改单；模型仍应结合调整费、实际库存和未来缺口决定是否修改订单。

\subsection*{9\quad 参数灵敏度与模型可信度}

准备量越大，紧急补购通常越少，但普通订单和未接收电量会增加，因此参数不能只按“越稳越好”来取。图\ref{fig:q3-parameter}在其余条件相同时比较了误差分位、日期衰减、分组长度和历史窗口。

在0.55--0.80的分位范围内，0.60的总费用最低；继续提高分位虽然压低了紧急费用，却使普通结算增加得更多。日期衰减取0.7时费用明显低于0.8和1.0，说明近期误差应当发挥更大作用。1、2、3、4和6小时分组中，3小时最低，且正好把每个6小时更新区间分为两段，在细节和稳定性之间取得平衡。历史窗口从14天改为28天或56天时，费用只相差十余元，因此选择能够覆盖四个星期的28天。

\begin{inlinefigure}
\includegraphics[width=.92\linewidth]{parameter_interpretability.pdf}
\caption{近期误差准备参数的费用依据}\label{fig:q3-parameter}
\end{inlinefigure}

此外，逐次撤去准备量、以及强制保留零点订单的25\%--100\%两类限制都会增加总费用，所以最终方案在四次决策中都保留完整准备量，也不给向下调整另设最低比例。这并不表示每次都必须改单，而是让式\eqref{eq:q3-objective}根据预报变化、调整费和真实库存自行判断。

为检查结果是否来自计算误差，对27条补充策略的1298592个实际10分钟区间和18980个准备量分组进行了独立回算。费用最大复算差约为 $1.9\times10^{-9}$ 元，库存、功率、订单锁定和逐段供需均满足约束。当前参数是在题目给出的2025年数据范围内比较得到，能够说明本年度的费用取舍和计算一致性；对于新的年份或不同小区，仍需用新数据重新校准，而不能直接认为同一组参数永远最优。
